%% ============================================================
%% sections/s3-amendment.tex
%% H. R. 9510 Bill v5.0 - SEC. 3. AMENDMENT TO THE FEDERAL FOOD, DRUG, AND
%% COSMETIC ACT. Carries the operative new section 515D (21 U.S.C. 360e-5)
%% verbatim from Bill v3.0 (subsections (a) through (j)), APPENDS the new v5.0
%% subsection (k) Financial Data Record, and carries the ten conforming
%% amendments, the clerical amendment, the rule of construction, and the
%% effective date, so the bill stands on its own.
%% Rendered visuals: Table 2 (ten-gate threshold schedule with unit costs),
%% Figure 3 (Mermaid BPMN decision rule with the cost-record checkpoint,
%% TikZ), Figure 4 (ASCII statutory layering with the financial sections),
%% Table 3 (ten conforming amendments crosswalk).
%% ============================================================
\billsec{SEC. 3.}{AMENDMENT TO THE FEDERAL FOOD, DRUG, AND COSMETIC ACT.}\phantomsection\label{toc:s3}

\uscprov{1}{\textbf{(a) In General.} Subchapter V of the Federal Food, Drug, and
Cosmetic Act (21 U.S.C. 351 et seq.) is amended by inserting after section 515C
(21 U.S.C. 360e-4) the following new section:}

\uscquote{0}{``SEC. 515D. VERIFICATION BEFORE GENERATION OF ROBOT-PATIENT
INTERACTION CODE IN PHYSICAL AI ONCOLOGY INVESTIGATIONS.}
\uscquote{0}{``(a) Requirement.- No robot-patient interaction code may be
generated or executed in a Physical AI oncology clinical investigation unless an
automated verification, validation, and uncertainty quantification process, bound
to one or more named external standards, has first cleared for that interaction,
and the cleared verification record has been documented and attested under
subsection (e).}
\uscquote{0}{``(b) Order of Operations.- A sponsor shall not generate, and shall
not execute, robot-patient interaction code in advance of the cleared
verification record. A record created after generation or execution does not
satisfy this section.}
\uscquote{0}{``(c) Gate Thresholds and Decision Rule.- For each robot-patient
interaction, the verification process shall apply the gate thresholds set out in
the schedule in Table 2. Each gate shall emit a determination of ACCEPT, BLOCK,
or ESCALATE; a single failing dimension shall set BLOCK, and ambiguity or
divergence shall set ESCALATE and return the interaction to a qualified human.
For each interaction the verification fraction shall equal 1.0, and, for an
interaction subject to a hard catastrophe predicate, that predicate shall also
hold \cite{asmevv40, nasastd7009}.}

% The threshold schedule is set at the body measure with left-aligned
% ragged-right columns; the v5.0 addition is the right-aligned illustrative
% unit-cost column, which feeds the per-gate schedule in Appendix B.
{\footnotesize
\begin{xltabular}{\textwidth}{@{}L{3.4cm} L{1.5cm} L{1.0cm} L{1.4cm} Y R{1.4cm}@{}}
\hline\noalign{\vskip 0.04cm}
\textbf{Gate} & \textbf{Min.\ agreement} & \textbf{Max.\ CV} & \textbf{Hard
predicate} & \textbf{Governing external standards} & \textbf{Cost/run (\$)}\\
\hline\noalign{\vskip 0.04cm}
\endfirsthead
\hline\noalign{\vskip 0.04cm}
\textbf{Gate} & \textbf{Min.\ agreement} & \textbf{Max.\ CV} & \textbf{Hard
predicate} & \textbf{Governing external standards} & \textbf{Cost/run (\$)}\\
\hline\noalign{\vskip 0.04cm}
\endhead
\hline\noalign{\vskip 0.04cm}
\endlastfoot
01 Bimanual hand-eye servo & 0.97 & 0.08 & No & IEC 80601-2-77; ISO 9283; ASME
V\&V 40 & 310\\
\hline\noalign{\vskip 0.04cm}
02 Dexterous finger force & 0.95 & 0.10 & No & IEC 80601-2-77; ISO/TS 15066 &
280\\
\hline\noalign{\vskip 0.04cm}
03 Whole-body balance & 0.98 & 0.06 & No & ISO 13482; IEC 60601-1 & 340\\
\hline\noalign{\vskip 0.04cm}
04 Auton.\ plan correct & 0.95 & 0.10 & No & UL 4600; IEEE 7009; IEC 62304 &
295\\
\hline\noalign{\vskip 0.04cm}
05 Instr.\ grasp, handover & 0.96 & 0.10 & No & IEC 80601-2-77; ISO 9283 & 270\\
\hline\noalign{\vskip 0.04cm}
06 Vasc.\ no-fly vol.\ (hand) & 1.00 & 0.05 & Yes & IEC 80601-2-77; ISO 14971 &
520\\
\hline\noalign{\vskip 0.04cm}
07 Bimanual suturing anast. & 0.96 & 0.08 & No & IEC 80601-2-77; Fistula Risk
Score & 305\\
\hline\noalign{\vskip 0.04cm}
08 Perception scene unders. & 0.95 & 0.10 & No & ASME V\&V 40; IEC 62304 & 290\\
\hline\noalign{\vskip 0.04cm}
09 Shared-room human coll. & 1.00 & 0.06 & Yes & ISO/TS 15066; ISO 10218-1; ISO
13482 & 560\\
\hline\noalign{\vskip 0.04cm}
10 Fault, emer.\ stop, degrad. & 1.00 & 0.05 & Yes & IEC 60601-1; ISO 13849-1;
IEEE 7009 & 540\\
\hline
\end{xltabular}}
\vspace{-0.8cm}
\figcaption{Table 2. Ten-gate threshold schedule with unit costs: minimum
validation agreement, maximum coefficient of variation, hard catastrophe
predicate, governing external standards, and the illustrative compute cost per
gate run (seed 20260525; the three hard-predicate gates carry added
catastrophe-simulation cost). The cost column sums to the 3,710 dollar
per-suite figure priced in Appendix B.}

\vspace{-0.2cm}

\uscprov{1}{The decision rule this schedule implements, with the new
financial-data checkpoint on the ACCEPT path, is shown in Figure 3.}

\begin{mermaidfig}
\node[circle,draw=billblack,fill=billwhite,minimum size=7mm,font=\scriptsize] (s) at (0,1.7) {Start};
\node[mmin]  (g)   at (0,0)       {Candidate robot-patient\\interaction code};
\node[mmdec] (q)   at (0,-2.4)    {Clears all 10\\standards-bound gates?};
\node[mmstep](esc) at (-4.3,-2.4) {ESCALATE\\human review};
\node[mmlaw] (cr)  at (4.3,-2.4)  {Cost record written\\per gate, to the audit\\trail (new 515D(k))};
\node[mmgoal](a)   at (4.3,-4.8)  {ACCEPT\\code may be generated};
\node[mmdec] (e)   at (0,-4.8)    {Divergence or\\ambiguity?};
\node[mmlaw] (b)   at (-4.3,-4.8) {BLOCK\\code withheld};
\draw[mmedge] (s) -- (g);
\draw[mmedge] (g) -- (q);
\draw[mmedge] (q) -- node[mmlabel]{yes} (cr);
\draw[mmedge] (cr) -- (a);
\draw[mmedge] (q) -- node[mmlabel]{no} (e);
\draw[mmedge] (e) -- node[mmlabel]{no} (b);
\draw[mmedge] (e) -- node[mmlabel,pos=0.35]{yes} (esc);
\draw[mmedge] (esc) -- (q);
\end{mermaidfig}
\vspace{-0.2cm}
\figcaption{Figure 3. The gate decision rule with the cost-record checkpoint
(gray-scale Mermaid as TikZ): ACCEPT, BLOCK, and ESCALATE keep their carried
semantics, and the v5.0 change is visible in the topology - no ACCEPT without
the 515D(k) cost record.}

\uscquote{0}{``(d) Readiness Gates.- Before a patient is enrolled or a robot
acts, the trial site shall pass the Physical AI Standard Level gate on
legislative authorization, regulatory compliance, and operational readiness;
each robot shall meet the minimum Unification Standard Level for its functional
type, which is not less than 7.0 for a surgical type, 6.0 for a therapeutic
positioning or diagnostic needle placement type, 4.0 for a rehabilitative type,
and 3.0 for a companion monitoring type; and each robot's behavior shall be
validated across not fewer than two simulation frameworks within the
cross-framework tolerances of less than 2 millimeters of trajectory deviation
and less than 0.5 newton of force discrepancy \cite{kawchak2026national}.}
\uscquote{0}{``(e) Documentation and Attestation.- Before robot-patient
interaction code is generated or executed, the sponsor shall make and file, for
each covered algorithm, algorithm documentation including a plain-language
summary, the decision logic or architecture, the training or conditioning data
sets, the gate bindings identifying the governing external standard and the
threshold enforced, the verification-before-generation record, and a
determinism statement giving the seed; a compliance statement and a signed
attestation by a responsible officer; and a hash-chained, tamper-evident audit
trail that satisfies part 11 of title 21, Code of Federal Regulations
\cite{cfr-part11, cfr-hipaa-164}.}
\uscquote{0}{``(f) Cybersecurity, Human Oversight, and Lifecycle.- Each Physical
AI system shall incorporate authentication and access control, encryption of
data in transit and at rest, network segmentation, and software-integrity
verification; shall operate under recorded human oversight with a manual
override and an emergency stop that brings the system to a safe state within
the procedure-appropriate budget and a hand-back-to-human escalation when a
gate escalates or a fault is detected; and shall maintain configuration and
model-version management, monitor for model and data drift, and, on
decommissioning, archive the audit trail, securely erase patient data, and
revoke access \cite{iec60601, iso13849}.}
\uscquote{0}{``(g) Nondiscrimination.- Each covered algorithm and its training
data shall minimize the risk of bias on protected characteristics and adhere to
evidence-based clinical guidelines, and each Physical AI system shall be
designed and operated consistent with applicable accessibility and
accommodation requirements \cite{cfr-92-210, usc-titlevi, usc-ada}.}
\uscquote{0}{``(h) Regulations.- Not later than 365 days after the date of
enactment of this section, the Secretary shall issue final regulations to carry
out this section. Pending those regulations, the Secretary may issue interim
guidance, and a sponsor may comply by meeting the requirements of subsections
(a) through (g).}
\uscquote{0}{``(i) Definitions.- In this section, the terms `Physical AI
system', `robot agent', `robot-patient interaction code', `generative
artificial intelligence', `AI-enabled device software function',
`verification', `validation', `uncertainty quantification', `verification
fraction', `validation agreement', `coefficient-of-variation bound', `hard
catastrophe predicate', `Physical AI Standard Level', `Unification Standard
Level', and `hash-chained audit trail' have the meanings given in the prior
bill and reconciled with section 201(h) and section 520(o) of this current
Act.}
\uscquote{0}{``(j) Rule of Construction.- Nothing in this section displaces the
investigational-device exemption under section 520(g), the protection of human
subjects under part 50 of title 21, Code of Federal Regulations, the
investigational new drug framework under part 312 of such title, or the
predetermined change control plan authority under section 515C as it applies
outside Physical AI, and nothing in this section itself reclassifies any
device.}
\uscquote{0}{``(k) Financial Data Record.-}
\uscquote{1}{``(1) Record.- For each verification run under this section, the
sponsor shall record the direct cost of the run - the compute consumed, the
records storage, and any qualified human review occasioned by an ESCALATE
determination - in the audit trail required by subsection (e), aggregated per
gate and per interaction class, in the controlled vocabulary of the Federal
budget process.}
\uscquote{1}{``(2) Attestation.- The verification cost record shall be covered
by the compliance statement and attestation under subsection (e), and shall be
maintained with the same hash-chained, tamper-evident integrity.}
\uscquote{1}{``(3) Annual Report.- The sponsor shall report the fiscal-year
aggregate of its verification cost records to the Secretary, for publication
in de-identified aggregate form under section 5 of the Verification Before
Generation in Physical AI Oncology Trials Act of 2026.}
\uscquote{1}{``(4) Definitions.- In this subsection, the term `verification
cost record' means the per-run financial data described in paragraph (1), and
the term `direct cost' means cost incurred specifically to execute, store, and
review a verification run, exclusive of general overhead.''.}

\uscprov{1}{\textbf{(b) Conforming Amendments.} The Federal Food, Drug, and
Cosmetic Act is amended, with the corresponding changes shown in the
comparative print in section 4, so that the device definition in section
201(h) reaches software that generates or executes robot-patient interaction
code; the prohibited acts in section 301 reach a violation of section 515D; a
device used without a cleared section 515D record is adulterated under section
501; the real world evidence program in section 505F recognizes section 515D
records; a verification-governed change consistent with a predetermined change
control plan needs no new premarket notification under section 510(k); special
controls under section 513 are graduated by degree of autonomy; a premarket
approval application under section 515 must include the cleared verification
record; a predetermined change control plan under section 515C must include an
automated verification-before-change protocol; the clinical-decision-support
exclusion in section 520(o) does not reach robot-control software; and a
savings clause at section 521 preserves State human-in-the-loop requirements
that are not different from, or in addition to, a Federal device requirement
\cite{usc-321h, usc-331, usc-351, usc-355g, usc-360, usc-360c, usc-360e,
usc-360e4, usc-360j, usc-360k}. Figure 4 shows where the new section attaches
on the device-regulation pathway and which financial sections it leans on, and
Table 3 is the conforming-amendment crosswalk.}

\begin{asciifig}
Robot-patient interaction code in a Physical AI oncology trial
      |
      v
Is it a "device"?  ................  sec. 321(h)            (s321)
      |   the sec. 360j(o) CDS exclusion does NOT carve out
      |   autonomous robot-control software ......  360j(o) (s360j)
      v
Classify ..........................  sec. 360c              (s360c)
Class II -> 510(k) notice .........  sec. 360(k)            (s360)
Class III -> premarket approval ...  sec. 360e              (s360e)
      +------> PCCP keystone ......  sec. 360e-4            (s360e-4)
               ==> NEW sec. 360e-5  (verify before generation
                   + the 515D(k) COST RECORD)

Financial sections the new rule leans on:
 - Device user fees (MDUFA V) .....  secs. 379i, 379j       (s379j)
 - Adulteration / quality .........  sec. 351               (s351)
 - Prohibited acts / enforcement ..  sec. 331               (s331)
 - Real-world evidence ............  sec. 355g              (s355g)
 - State preemption boundary ......  sec. 360k              (s360k)
\end{asciifig}
\vspace{-0.6cm}
\figcaption{Figure 4. Statutory layering with the financial sections (ASCII):
the device question, the carve-out the code must not fall through,
classification, the premarket pathways, the section 360e-4 keystone, the new
section 360e-5 with its cost record, and the financial sections - including
the user-fee sections - the rule leans on.}

{\footnotesize
\begin{xltabular}{\textwidth}{@{}L{1.0cm} L{1.6cm} L{1.7cm} Y@{}}
\hline\noalign{\vskip 0.04cm}
\textbf{No.} & \textbf{FD\&C \S} & \textbf{21 U.S.C.} & \textbf{Conforming change
made by H. R. 9510}\\
\hline\noalign{\vskip 0.04cm}
\endfirsthead
\hline\noalign{\vskip 0.04cm}
\textbf{No.} & \textbf{FD\&C \S} & \textbf{21 U.S.C.} & \textbf{Conforming change
made by H. R. 9510}\\
\hline\noalign{\vskip 0.04cm}
\endhead
\hline\noalign{\vskip 0.04cm}
\endlastfoot
(1) & \S~201(h) & \S~321(h) & Device definition reaches software that generates
or executes robot-patient code\\
\hline\noalign{\vskip 0.04cm}
(2) & \S~301 & \S~331 & New prohibited act for generation or execution in
violation of section 515D\\
\hline\noalign{\vskip 0.04cm}
(3) & \S~501 & \S~351 & Device used without a cleared section 515D record is
adulterated\\
\hline\noalign{\vskip 0.04cm}
(4) & \S~505F & \S~355g & Real world evidence program recognizes section 515D
records\\
\hline\noalign{\vskip 0.04cm}
(5) & \S~510(k) & \S~360(k) & Verification-governed change needs no new
premarket notification\\
\hline\noalign{\vskip 0.04cm}
(6) & \S~513 & \S~360c & Special controls graduated by degree of autonomy\\
\hline\noalign{\vskip 0.04cm}
(7) & \S~515 & \S~360e & Premarket approval application must include the
cleared verification record\\
\hline\noalign{\vskip 0.04cm}
(8) & \S~515C & \S~360e-4 & Change control plan must include automated
verification-before-change protocol\\
\hline\noalign{\vskip 0.04cm}
(9) & \S~520(o) & \S~360j(o) & Clinical-decision-support exclusion does not
reach robot-control software\\
\hline\noalign{\vskip 0.04cm}
(10) & \S~521 & \S~360k & Savings clause preserves State human-in-the-loop
requirements\\
\hline
\end{xltabular}}
\vspace{-0.8cm}
\figcaption{Table 3. The ten conforming amendments crosswalk, carried from the
prior versions: each amendment by its Federal Food, Drug, and Cosmetic Act
section, its 21 U.S.C. citation, and the change made. The eleventh, financial
amendment - the user-fee treatment in \S~738 - is made by SEC. 5(c) and shown
in the section 4 print.}

\uscprov{1}{\textbf{(c) Clerical Amendment.} The table of contents for the
Federal Food, Drug, and Cosmetic Act is amended by inserting after the item
relating to section 515C the following new item:}
\uscquote{0}{``Sec. 515D. Verification before generation of robot-patient
interaction code in Physical AI oncology investigations.''.}

\uscprov{1}{\textbf{(d) Rule of Construction.} Nothing in this Act displaces
the investigational-device exemption under section 520(g) of the Federal Food,
Drug, and Cosmetic Act (21 U.S.C. 360j(g)), the protection of human subjects
under part 50 of title 21, Code of Federal Regulations, the investigational
new drug framework under part 312 of such title, the predetermined change
control plan authority under section 515C as it applies outside Physical AI,
or any State law preserved under section 521 as amended; and nothing in this
Act reclassifies any device \cite{cfr-part50, cfr-part312}.}

\uscprov{1}{\textbf{(e) Effective Date; Implementation Deadline; Transition
Rule.} This Act takes effect on the first day of the first calendar quarter
beginning not less than one year after enactment; the Secretary shall issue
the final regulations required by section 515D(h) not later than 365 days
after enactment; a Physical AI oncology clinical investigation already
enrolling on the effective date shall have 180 days to come into compliance
with section 515D; the first annual financial-data report under section
515D(k)(3) is due for the first full fiscal year beginning after the effective
date; and this Act does not disturb the February 2, 2026, effective date of
the Quality Management System Regulation \cite{fr-qmsr-2024}.}
